\documentclass[a4paper,12pt]{article}

\usepackage[catalan]{babel}
\usepackage{fontspec}
\usepackage{fullpage}
\usepackage[a4paper, margin=2cm]{geometry} % To change the margins
\usepackage{graphicx} % Insert images
\usepackage[hidelinks]{hyperref} % Links color
\usepackage{import}
\usepackage{ragged2e}
\usepackage{wrapfig} %To Text wrap
\usepackage{listings} % Add code
\usepackage{tikz}
\usepackage{calc}


%Tikz defines
\usetikzlibrary{shapes, arrows, positioning, automata, calc}
\tikzstyle{UML} = [rectangle split, rectangle split parts=3, draw,text centered,
				   every two node part/.style={align=left, font=\scriptsize}, 
				   every three node part/.style={align=left, font=\scriptsize},
				   font=\scriptsize\bf]
\tikzstyle{cas} = [ellipse, draw, text centered, font=\scriptsize]
\tikzstyle{arrow} = [thick,-]
\begin{document}
\title{\textsc{Ken-Ken} \\ \large 1a entrega PROP}
\author{Marc Asenjo i Ponce de León \and
	Arnau Canyadell i Miquel \and
	Joan Marcè i Igual \and
	Iñigo Moreno i Caireta \and
	Esteve Tarragó i Sanchís}

\date{\today}
\maketitle

El nostre grup ha escollit fer el treball de PROP sobre el Ken-Ken. (explicació
sobre el Ken-Ken?). 

En aquest cas l'usuari podrà gestionar Ken-Kens generats per ell o també podrà
deixar que la maquina li generi un. 
Per als Ken-Ken que l'usuari crei el programa comprovarà que hi hagi una solució
possible sinó, el donara com a \emph{no vàlid} i no el deixarà introduir. Els Ken-Ken generats 
pel programa podran tenir uns patrons predefinits tals com la mida del tauler, la forma de les 
regions (hi haurà una llista on es podran escollir les que es volen i les que no) o el nombre 
de nombres ja co\lgem ocats.

En tots dos casos l'usuari podrà intentar resoldre el Ken-Ken desitjat amb
l'assistència del programa que anirà comprovant el compliment de les regles del
joc, en el cas que no es complís alguna norma es mostrarà un avís. Durant el
desenvolupament del joc es tindran en compte les accions de l'usuari per
atorgar-li una puntuació en funció del temps que hagi durat la partida
i també dels errors comesos per l'usuari durant aquesta.

A part també hi haurà un rànquing on quedaran registrades les millors puntuacions,
hi haurà el nom de l'usuari i la puntuació que aquest ha fet.

\section{Casos d'ús}

\subsection{Diagrama}
\begin{center}
\begin{tikzpicture}[node distance=0.5cm and 0.6cm, auto, align=center]
\node(StartApp){\includegraphics[width=1.3cm]{./Images/stickman}};
\node[cas, right = of StartApp](User){Entrar\\usuari};
\node[cas, right = of User](Ranking){Veure\\rànquing};
\node[cas, above = of Ranking, yshift=5.5cm](Start){Començar\\a jugar};
\node[cas, below = of Ranking](Admin){Veure co\lgem ecció\\de Ken-Kens};
\node[cas, below = of User](Exit){Sortir};

%Començar a jugar
\node[cas, right = of Start](Database){Seleccionar\\tauler de la\\co\lgem ecció};
\node[cas, above = of Database, xshift=1cm](Load){Carregar\\joc};
\node[cas, below = of Database, xshift=-1cm](CreateCPU){Generar i\\jugar};

%Rànquing
\node[cas, right = of Ranking](RRanking){Netejar\\dades};

%Veure col·lecció
\node[cas, right = of Admin](SelectKK){Seleccionar\\Ken-Ken};
\node[cas, below = of Admin](Create){Crear\\Ken-Ken};

\node[cas, right = of CreateCPU, xshift=0.5cm](Play){Jugar};

%Generar i jugar
\node[cas, below = of CreateCPU](SelectSize){Seleccionar\\Mida};
\node[cas, below = of SelectSize](SelectRegions){Seleccionar núm.\\ i mida regions};

%Jugar
\node[cas, right = of Play](Hint){Pista};
\node[cas, above = of Hint](SelectCel){Seleccionar\\ce\lgem a};
\node[cas, below = of Hint](Save){Guardar};
\node[cas, below = of Save](Surrender){Rendir-se};

%Seleccionar cel·la
\node[cas, above right = of SelectCel](ChangeNum){Canviar\\nombre};
\node[cas, right = of SelectCel](DeleteNum){Eliminar\\nombre};
\node[cas, below right = of SelectCel](Notify){Notificació};

%Seleccionar Ken-Ken
\node[cas, right = of SelectKK](EditKK){Editar\\Ken-Ken};
\node[cas, above right = of SelectKK](DeleteKK){Eliminar\\Ken-Ken};
\node[cas, below right = of SelectKK](ViewSolution){Veure\\Solució};

%Crear Ken-Ken
\node[cas, below left = of Create, yshift=-3.5cm](ChooseSize){Triar\\mida};

%Triar mida
\node[cas, right = of ChooseSize, xshift=1cm](CreateRegion){Crear Regió};
\node[cas, above = of CreateRegion](SelectCell){Seleccionar\\ce\lgem a};
\node[cas, above = of SelectCell](FillRandom){Emplenar de\\nombres aleatoris\\correctes};
\node[cas, below = of CreateRegion, xshift=1cm](SelectRegion){Seleccionar\\regió};
\node[cas, below = of SelectRegion, xshift=1cm](ViewSolutionable){Veure si\\té solució};
\node[cas, below = of ViewSolutionable](ViewUniqueSolution){Veure si té\\solució única};
\node[cas, left = of ViewUniqueSolution](SaveNew){Guardar};
\node[cas, below = of ChooseSize, yshift=-2cm](Solve){Solucionar};

%Seleccionar cel·la
\node[cas, right = of SelectCell, yshift=1cm, xshift=2cm](SetNum){Anotar nombre\\correcte};
\node[cas, above= of SetNum, yshift=-0.5cm](RandomNum){Nombre aleatori\\correcte};

%Crear regió
\node[cas, right = of CreateRegion, yshift=1cm, xshift=2cm](ChangeNum2){Canviar nombre\\ i/o resultat};
\node[cas, below = of ChangeNum2, yshift=0.5cm](SelectCells){Seleccionar\\ce\lgem es};

%Seleccionar regió\\
\node[cas, right = of SelectRegion, xshift=2.2cm](ChangeCells){Modificar\\ce\lgem es};
\node[cas, below = of ChangeCells, yshift=0.5cm](ChangeNumRes){Modificar nombre\\i/o resultat};

\foreach \x/\y in {User/Start, User/Ranking, User/Admin, Start/Database, Start/CreateCPU,
				   Ranking/RRanking, Admin/SelectKK, Admin/Create, Database/Play, 
				   Play/SelectCel, Play/Surrender.north west, Play/Hint, SelectCel/ChangeNum, 
				   SelectCel/DeleteNum, SelectCel/Notify, Play/Save, SelectKK/Play, Start/Load.west, 
				   Load.south east/Play, User/Exit, CreateCPU/SelectSize, SelectKK/ViewSolution,
				   SelectSize/SelectRegions, SelectRegions.north east/Play, StartApp/User,
				   SelectKK/DeleteKK, SelectKK/EditKK, Create/ChooseSize, 
				   ChooseSize/FillRandom.south west, ChooseSize/CreateRegion, ChooseSize/SelectCell, 
				   ChooseSize/SelectRegion, ChooseSize/ViewSolutionable, ChooseSize/ViewUniqueSolution, 
				   ChooseSize/Solve, SelectCell/RandomNum.south west, SelectCell/SetNum.west, 
				   CreateRegion/ChangeNum2.west, CreateRegion/SelectCells.west, 
				   SelectRegion/ChangeCells, SelectRegion/ChangeNumRes, ChooseSize/SaveNew} 
{
\draw[thick, ->] (\x) -- (\y);
};

\end{tikzpicture}
\end{center}

\subsection{Descripció}


\paragraph{Nom:} Entrar Usuari \\%Set name here
\textbf{Descripció:}
Permet introduir el nom d'usuari perquè es comprovi a la base de dades si ja hi ha un usuari i sinó crear-ne un.

\paragraph{Nom:}  Començar a jugar \\
\textbf{Descripció:} 
Obre la pantalla de selecció dels Ken-Ken per carregar un joc anterior, generar un nou tauler o seleccionar-ne un de ja creat.

\paragraph{Nom:} Carregar joc \\
\textbf{Descripció:}
Permet continuar jugant una partida guardada anteriorment.

\paragraph{Nom:} Seleccionar tauler de la co\lgem ecció \\
\textbf{Descripció:}
Començar una nova partida utilitzant un tauler dels que hi ha a la co\lgem ecció de taulers.

\paragraph{Nom:} Generar i jugar \\
\textbf{Descripció:}
El programa genera un nou Ken-Ken a partir dels paràmetres especificats de l'usuari després el podrà començar a jugar. Farà avançar el menú per poder escollir la mida del tauler.

\paragraph{Nom:} Seleccionar mida \\
\textbf{Descripció:}
Permet escollir la mida que s'utilitzarà per generar un nou Ken-Ken ($3x3, 4x4, 5x5... 9x9$).
Permetrà continuar el menú per poder escollir el número i mida de les regions.

\paragraph{Nom:} Seleccionar número i mida regions \\
\textbf{Descripció:} 
Permetrà escollir de quina mida seran les regions del Ken-Ken generat ($1 - 5$) i després ja passarà a crear el tauler i que l'usuari el comenci a jugar.

\paragraph{Nom:} Veure rànquing\\
\textbf{Descripció:} 
Permet veure el rànquing amb els millors usuaris i les seves puntuacions i estadístiques.

\paragraph{Nom:} Jugar\\
\textbf{Descripció:} 
Permet veure el tauler actual, amb els nombres i possibilitats que ja estan co\lgem ocades, seleccionar una ce\lgem a, i accedir a altres opcions.

\paragraph{Nom:} Seleccionar ce\lgem a\\
\textbf{Descripció:} 
La ce\lgem a seleccionada queda destacada i es podrà escollir que fer amb ella: canviar el seu valor, deixar-la sense o afegir possibilitats.


\paragraph{Nom:} Canviar nombre\\
\textbf{Descripció:} 
Queda marcat el valor de la ce\lgem a al tauler.


\paragraph{Nom:} Eliminar nombre\\
\textbf{Descripció:} 
Queda eliminat el valor de la ce\lgem a.

\paragraph{Nom:} Notificació\\ %CANVIAR NOM?
\textbf{Descripció:} 
Permet anotar la possibilitat que un nombre pugui estar en la ce\lgem a seleccionada.

\paragraph{Nom:} Pista\\
\textbf{Descripció:} 
Es resol una ce\lgem a automàticament.


\paragraph{Nom:} Guardar\\
\textbf{Descripció:} 
Es guarda el progrés de la partida en la co\lgem ecció de Ken-Kens


\paragraph{Nom:} Rendir-se\\
\textbf{Descripció:} 
Permet abandonar la partida i tornar a l'inici de l'aplicació.



%Add descripcion here

\newpage
\section{Diagrama de classes}
%ADD ARROWS
\begin{center}
\begin{tikzpicture}[node distance=1cm and 1.3cm, auto, align=center]
\node[UML](Application){Application \nodepart{two} };
\node[UML, below = of Application](KanKun){KanKun};
\node[UML, below = of KanKun](GameLauncher){GameLauncher};

\node[UML, right = of Application, xshift=0.6cm](DB){DB};
\node[UML, right = of DB](Table){Table};
\node[UML, below = of Table](Data){Data \nodepart{three} +save()\\+load()};
\node[UML, right = of Table](Ranking){Ranking};
\node[UML, right = of Data](Game){Game \nodepart{two}
	IdGame\\
	Score\\
	Status};
\node[UML, below = of Game](Player){Player \nodepart{two} Nom};

\node[UML, below left = of Data](Board){Board};
\node[UML, below = of Board, yshift = -0.5cm](KKBoard){KKBoard};
\node[UML, right = of Board](Cell){Cell};
\node[UML, right = of KKBoard](Region){Region};

\node[UML, below = of KKBoard](BoardCreator){BoardCreator};
\node[UML, below right = of BoardCreator](BoardHumanEditor){BoardHumanEditor};
\node[UML, below = of BoardCreator](BoardCPUGenerator){BoardCPUGenerator};

\node[left = of GameLauncher, xshift = 1cm](A){};


\draw [arrow, -latex'] (KanKun) -- (Application);
\draw [arrow, dashed, -latex'] (GameLauncher) -- (KanKun);
\draw [arrow, dashed, -latex'] ($(DB.south west)+(0, 0.4)$) -- ++ (-1, 0) |- (GameLauncher.east);
\draw [arrow, dashed, -latex'] (DB) -- (Application)
	node[font=\scriptsize, above right, xshift=1cm]{<<static>>};
\draw [arrow, -diamond] (Table.west) -- (DB.east);
\draw [arrow, -latex'] (Ranking.west) -- (Table.east);
\draw [arrow, -diamond] (Data) -- (Table);
\draw [arrow, -latex'] (Game) -- (Data);
\draw [arrow, -diamond] (Game.north) -- ++ (0, 0) -| (Ranking.south);
\draw [arrow] (Player.west) -- ++ (-0.5, 0) |- (Data.east);
\draw [arrow] (Player) -- (Game)
	node[font=\scriptsize, above right, at start]{1}
	node[font=\scriptsize, below right, at end]{*};
\draw [arrow, -latex'] (Board) |- (Data);
\draw [arrow, -diamond] (Cell) -- (Board);
\draw [arrow, -latex'] (KKBoard) -- (Board);
\draw [arrow, -diamond] (Region) -- (KKBoard);
\draw [arrow] (Region.north) --  ++(0, 0.5) -| (Cell) 
	node[font=\scriptsize, above right, at start, yshift=-0.5cm]{1}
	node[font=\scriptsize, below right, at end]{*};
\draw [arrow, dashed, -latex'] (KKBoard) -- (BoardCreator);
\draw [arrow, -latex'] (BoardHumanEditor) |- (BoardCreator);
\draw [arrow, -latex'] (BoardCPUGenerator) -- (BoardCreator);
\draw [arrow, -latex', dashed] (BoardHumanEditor.south) -- ++ (0, -0.5) -| (A.center) |- (KanKun.west);
\draw [arrow, -latex', dashed] (BoardCPUGenerator) -| (GameLauncher);
\end{tikzpicture}
\end{center}
\end{document}